\subsection{The final transfer graph}\label{sec:transfer-graph}

This section makes precise how the two criteria relate once every selected
pair has been identified. We show that, from that point on, the orthogonality
phase of $\SG$ is exact: it fails exactly when the displayed forms satisfy a
nontrivial quadratic relation. We use only simplicity: no cell lies in two
displayed squares. The base need not be $C_4$-free or maximum. The rules of
$\RW$ and $\SG$, and their soundness proofs in \cite{LQ}, use only
simplicity, although they are stated there for limited graphs.

\subsubsection{The graph and its relations}

Let $(F,D)$ be a simple display. Its \emph{classes} are the displayed
squares' supports: the singletons of $F$ and the pairs of $D$. For a
class $X$ put $f_X=\sum_{a\in X}u_a$. Thus
\[
 P_{F,D}=\sum_X f_X^2 \qquad\text{and}\qquad N=|F|+|D|.
\]

\begin{definition}\label{def:final-graph}
The \emph{final transfer graph} $\Gamma(D)$ is defined as follows.
\begin{itemize}
\item \emph{Vertices} are the unordered pairs $\{X,Y\}$ of distinct
      classes. A \emph{cross pair} of $\{X,Y\}$ is a pair of cells
      $x\in X$, $y\in Y$.
\item A vertex is \emph{grounded} if some cross pair
      \begin{enumerate}[label=(\roman*)]
      \item lies in one row or one column, or
      \item spans a genuine rectangle whose other diagonal contains a
            hole or is a pair of $D$.
      \end{enumerate}
\item \emph{Edges}: each genuine rectangle with all four corners occupied,
      whose two diagonals are cross pairs of vertices $v$ and $w$, gives an
      edge $vw$. It is a loop when $v=w$.
\end{itemize}
A component is \emph{grounded} if it contains a grounded vertex. It is
\emph{bipartite} if it has no closed walk of odd length.
\end{definition}

Every loop is at a grounded vertex. Suppose the two diagonals
$\{(r,c),(s,d)\}$ and $\{(r,d),(s,c)\}$ of one rectangle give the same
vertex. Then two corners in one row, or in one column, belong to the same
class. Another cross pair of that vertex then lies in one line, so the
vertex is grounded by~(i).

\Needspace{10\baselineskip}
\begin{theorem}\label{thm:final-graph-relations}
Let $C=(C_{XY})$ be a real symmetric matrix indexed by classes. Then
$\sum_{X,Y}C_{XY}f_Xf_Y$ vanishes identically if and only if $C$ has zero
diagonal and the function $h(\{X,Y\})=C_{XY}$ satisfies:
\begin{itemize}
\item $h$ vanishes at grounded vertices;
\item $h(v)=-h(w)$ along every edge $vw$.
\end{itemize}
The space of such relations therefore has dimension equal to the number of
ungrounded bipartite components of $\Gamma(D)$.

If such a component exists, then $\sos(P_{F,D})\le N-1$.
\end{theorem}

\begin{proof}
We compare coefficients of the biquadratic monomials.
\begin{itemize}
\item \emph{Squares $u_a^2$.} This monomial occurs only in $f_X^2$ for the
      class $X$ of $a$, with coefficient $C_{XX}$. Hence $C_{XX}=0$.
\item \emph{Line monomials} $x_r^2y_cy_d$ and $x_rx_sy_c^2$. Each arises
      from a single cell pair, and when its cells lie in distinct classes
      the coefficient is $2h(v)$. This is condition~(i).
\item \emph{Rectangle monomials} $x_rx_sy_cy_d$. Each arises exactly from
      the two diagonals of one rectangle. A diagonal contributes $2h$ of its
      vertex when both cells are occupied and in distinct classes, and
      nothing otherwise (a hole, or a single class, where $C_{XX}=0$). So
      the coefficient vanishes exactly when $h(v)+h(w)=0$ along an edge,
      or $h(v)=0$ under ground~(ii).
\end{itemize}
On a component, these conditions force $h\equiv0$ unless the component is
ungrounded and bipartite. On such a component the solutions are the
multiples of its $\pm1$ colouring.

Now take such a colouring $H$. It is symmetric, nonzero, has zero diagonal
and trace, and satisfies $f^{\mathsf T}Hf\equiv0$. With
$\lambda=\lambda_{\max}(H)>0$, the matrix $W=I-H/\lambda$ is positive
semidefinite of rank at most $N-1$, and $P_{F,D}=f^{\mathsf T}Wf$.
Factoring $W$ gives at most $N-1$ bilinear squares; this is the rank step
of \cite[Lemma~2.13]{XY}.
\end{proof}

\subsubsection{Completeness after identification}

In any SOS representation of $P_{F,D}$, the coefficient equations of
Section~\ref{sec:prelim} hold for the cell vectors. Every rule of $\RW$
and $\SG$ is a deduction from these equations, together with the
Cauchy--Schwarz inequality for the unit vectors of a selected pair. We say that a criterion
\emph{identifies} $D$ if its closure resolves every selected pair.

\begin{theorem}\label{thm:orthogonality-completeness}
Let $(F,D)$ be simple.
\begin{enumerate}[label=(\alph*)]
\item If $(F,D)$ satisfies $\SG$, then $\Gamma(D)$ has no ungrounded
      bipartite component.
\item If $\SG$ identifies $D$ and $\Gamma(D)$ has no ungrounded bipartite
      component, then $(F,D)$ satisfies $\SG$.
\item If $(F,D)$ satisfies $\RW$, then every component of $\Gamma(D)$ is
      grounded. Conversely, if $\RW$ identifies $D$ and every component is
      grounded, then $(F,D)$ satisfies $\RW$.
\end{enumerate}
In particular, if $\SG$ identifies $D$, the following are equivalent:
\begin{itemize}
\item $(F,D)$ satisfies $\SG$;
\item $\Gamma(D)$ has no ungrounded bipartite component;
\item $P_{F,D}$ is irreducible.
\end{itemize}
\end{theorem}

\begin{proof}
\emph{(a).} Let $K$ be an ungrounded bipartite component with colouring
$H$, and take $0<t\le1/\lambda_{\max}(H)$. Then $I-tH$ is positive
semidefinite with unit diagonal; write $I-tH=B^{\mathsf T}B$ with columns
$w_X$. Put $v_a=w_X$ for each $a\in X$, and $v_h=0$ at holes. Since
$f^{\mathsf T}Hf\equiv0$, these vectors come from the SOS representation
\[
 P_{F,D}=\sum_i\bigl(\textstyle\sum_a(v_a)_iu_a\bigr)^2 .
\]
For $x\in X$ and $y\in Y$ with $\{X,Y\}\in K$, we get
$\langle v_x,v_y\rangle=-tH_{XY}\ne0$. But an accepting closure
certifies $\langle v_x,v_y\rangle=0$ in every representation.

\emph{(b).} Once every pair is resolved, the current classes are the
displayed supports. The signed transfer graph of that stage then contains
all grounds and edges of $\Gamma(D)$:
\begin{itemize}
\item the line rule gives ground~(i);
\item transfer from a hole, or from an identified pair's diagonal, gives
      ground~(ii);
\item the rectangle identities give the sign-reversing edges, since all
      prescriptions between distinct final classes vanish.
\end{itemize}
Grounded components are forced by propagation, and every other component
has an odd cycle. So every pair of classes is certified orthogonal.

\emph{(c).} Every certified orthogonality $p\perp q$, with $p,q$ in
distinct final classes $X,Y$, lies in a grounded component of $\Gamma(D)$.
We prove this by induction over the $\RW$ derivation.
\begin{itemize}
\item \emph{Line rule.} This gives ground~(i).
\item \emph{Saturation.} This stays within the same vertex, because
      identifications merge only the two cells of a selected pair.
\item \emph{Complementary-pair rule.} This only identifies pairs and
      produces no orthogonality.
\item \emph{Rectangle transfer} from the other diagonal $\{r,s\}$.
      \begin{itemize}
      \item If that diagonal contains a hole, or is an identified pair, we
            get ground~(ii).
      \item If $r\perp s$ was certified, then $\delta(\{r,s\})=0$, so
            $\{r,s\}$ is not a selected pair and $r$, $s$ lie in distinct
            classes. The rectangle is an edge to their vertex, which lies in
            a grounded component by induction.
      \end{itemize}
\end{itemize}
Acceptance certifies every vertex. The converse is as in~(b), without
odd cycles. Saturation moves a certificate between the cross pairs of one
vertex, and rectangle transfer carries it along an edge.

For the final equivalence: (a) and (b) give the first equivalence;
Theorem~\ref{thm:final-graph-relations} and the soundness of $\SG$
\cite[Theorem~7.2]{LQ} give the second.
\end{proof}

\begin{remark}\label{rem:final-graph}
\emph{The identification hypothesis is necessary.} Consider the
$4\times4$ display $T$ of \cite[(40)]{XY}: the four diagonal cells as
singletons, plus the six pairs $F^+$ of \cite[(38)]{XY}.
\begin{itemize}
\item A direct check shows that all $15$ components of $\Gamma$ are
      grounded. By Theorem~\ref{thm:final-graph-relations} there is no
      product relation, in agreement with \cite[Proposition~5.4]{XY}.
\item The display has no holes, no degenerate pairs and no complementary
      pairs, so neither closure identifies any pair.
\item Yet $T$ is a sum of nine squares \cite[Theorem~5.3]{XY}.
\end{itemize}
As Xu and Yan remark after their Proposition~5.4, shorter representations
may use forms outside the displayed span. Identification is what excludes
them here.

\emph{Related work.} The proof of \cite[Proposition~5.4]{XY} already
observes that relations have zero diagonal and correspond to the kernel of
the product map for one display. \cite[Remark~7.8]{LQ} describes the Gram
perturbation mechanism. In their formulation,
\cite[Proposition~2.3]{CLQIncidence} shows that after an identification
prefix, grounding of every component suffices for $\RW$; this is the
converse half of Theorem~\ref{thm:orthogonality-completeness}(c).
Here we add the necessity direction, the description of all product
relations (Theorem~\ref{thm:final-graph-relations}), and the equivalence
of $\SG$ with irreducibility when the signed closure identifies $D$.

\emph{Relation to \cite[Remark~7.8]{LQ}.} That remark declines a converse
for ungrounded bipartite signed components. The remark remains appropriate
whenever identification is incomplete, including at fixed points with
unresolved pairs. Theorem~\ref{thm:final-graph-relations} gives the
converse at the final round.

\emph{Consequence for the comparison.} If $\RW$ identifies $D$, then
$(F,D)$ satisfies $\SG$ but not $\RW$ exactly when $\Gamma(D)$ has an
ungrounded odd component and no ungrounded bipartite one. The $15\times6$
witness of \cite[Theorem~7.5]{LQ} and the display of
Proposition~\ref{prop:full-four-distinction} are of this kind; the latter's
three odd components have sizes $5$, $13$ and $3$. In any other
signed-only display, $\RW$ fails to identify some pair.

\emph{Before any identification.} Consider a hole-free display. Before
any identification, every current class has one cell, so each vertex has a single
cross pair and at most one rectangle. The only pairs either rule can then
resolve are degenerate and complementary pairs, and the two closures agree.

\end{remark}
